%%%%%%%%%%%%%%%%%%%%%%%%%%%%%%%%%%%%%%%%%%%%%%%%%%%%%%%%%%%%%%%%%%%%%%%%%%%%%%%
%% ACM MobiCom 2027 submission --- ActDelta
%%
%% !!! IMPORTANT !!!
%% Every quantitative result in this draft is a PLACEHOLDER consistent with the
%% hypotheses and the pre-registered pass/fail thresholds in the project
%% specification (IDEA2-closed-loop-predictive-transmission.md).  They are NOT
%% measured numbers.  Replace every one of them before submission.
%%%%%%%%%%%%%%%%%%%%%%%%%%%%%%%%%%%%%%%%%%%%%%%%%%%%%%%%%%%%%%%%%%%%%%%%%%%%%%%
\documentclass[sigconf,10pt,anonymous,review]{acmart}

%% ---------------------------------------------------------------------------
%% Submission-time settings (revert for camera-ready)
%%
%% The ACM Reference Format / permission / ISBN / DOI block on page 1 is left
%% at the acmart default (printed).  Fifteen of the twenty-one MobiCom papers
%% we measured carry it, acmart prints it unless told not to, and it is
%% mandatory at camera-ready --- so `printacmref=false' only hides a problem.
%% In anonymous mode acmart suppresses the author names inside it for us.
%% ---------------------------------------------------------------------------
\setcopyright{acmlicensed}
\settopmatter{printfolios=true}

%% The CFP forbids embedded hyperlinks in the submitted PDF.
%% Uncomment before generating the PDF you upload to HotCRP.
% \hypersetup{draft}

%% SIGMOBILE had not announced the 2027 venue as of 2026-08-19 (the site lists
%% the dates only).  Replace the third argument with the exact string the
%% organizers publish before camera-ready.
\acmConference[ACM MobiCom '27]{The 33rd Annual International Conference on
  Mobile Computing and Networking}{October 18--22, 2027}{To be announced}
\acmYear{2027}
\copyrightyear{2027}
\acmISBN{978-1-4503-XXXX-X/2027/10}
\acmDOI{10.1145/nnnnnnn.nnnnnnn}

%% ---------------------------------------------------------------------------
%% Packages
%% ---------------------------------------------------------------------------
\usepackage{booktabs}
\usepackage{multirow}
\usepackage{subcaption}
\usepackage{amsmath}
\usepackage{needspace}

%% ---------------------------------------------------------------------------
%% Run-in paragraph headings.  acmart/sigconf sets them in italic
%% (\def\@parfont{\itshape}); every MobiCom paper we measured sets them bold,
%% and reserves italic for emphasis.  Match the venue, not the class default.
%% ---------------------------------------------------------------------------
\makeatletter
\def\@parfont{\bfseries}
\makeatother

%% ---------------------------------------------------------------------------
%% Figures
%%
%% Every figure is compiled separately and included here as a PDF.  See
%% figures/README.md; the short version is that each figure is three files with
%% the same stem --- <fig>.tex (drawing code), <fig>.json (data), <fig>.pdf ---
%% and `make` in figures/ rebuilds the PDFs.  This document therefore loads no
%% TikZ and no pgfplots at all.
%%
%% The PDFs are included at their natural size, never scaled: a standalone
%% figure is built against acmart's own \columnwidth and font sizes, so scaling
%% it here would be the one thing that makes figure text stop matching body
%% text.  If a figure is the wrong size, change its width in its own .tex.
%% ---------------------------------------------------------------------------
\newcommand{\figexp}[1]{\includegraphics{figures/experiment/#1}}
\newcommand{\figsch}[1]{\includegraphics{figures/schematic/#1}}

%% Shared with every standalone figure, so figure text and body text cannot
%% drift.  Holds \sys and friends, the colour palette, and \cnum --- the latter
%% two because \cnum is used in Figure 1's caption, which lives here.
%% ---------------------------------------------------------------------------
%% Everything main.tex AND the standalone figures both need.
%%
%% \input by main.tex and by figures/preamble.tex, so a macro defined here
%% renders identically in the body text, in a caption, and inside a figure.
%% Never define these anywhere else --- a second definition is how figure text
%% and body text drift apart.
%%
%% The rule for what lives here rather than in preamble.tex: if main.tex refers
%% to it, it belongs here.  \cnum is the reason the colours are here too --- it
%% is used both inside Figure 1 and in that figure's caption, and captions stay
%% in main.tex.
%% ---------------------------------------------------------------------------
\usepackage{xspace}

\newcommand{\sys}{ActDelta\xspace}
\newcommand{\pizero}{$\pi_0$\xspace}
\newcommand{\pizerofive}{$\pi_{0.5}$\xspace}
\newcommand{\pp}{\,pp\xspace}

\definecolor{cA}{RGB}{31,119,180}
\definecolor{cB}{RGB}{214,39,40}
\definecolor{cC}{RGB}{44,160,44}
\definecolor{cD}{RGB}{255,127,14}
\definecolor{cE}{RGB}{148,103,189}
\definecolor{cGrey}{RGB}{120,120,120}

%% circled step number, safe inside captions (no tikz in moving arguments)
\newcommand{\cnum}[1]{\textcolor{cB}{\textcircled{\raisebox{0.35pt}{\tiny #1}}}}


\begin{document}

%% The optional argument is the running head.  Do not put \sys there: the macro
%% ends in \xspace, which inserts a space before the colon when the title is
%% re-typeset in the header.  A genuine short title also keeps the header from
%% running the full width of the page, which is what the MobiCom corpus does.
\title[ActDelta: Decision-Triggered Predictive Transmission]{\sys:
  Decision-Triggered Predictive Transmission for Edge-Offloaded Robot Policies}

\author{Anonymous Author(s)}

\begin{abstract}
A robot that cannot host a $7$B vision-language-action (VLA) policy uploads its
cameras to an edge server and receives action chunks back. The uplink is the
bottleneck: the $11$\,Mbps a deployed policy server asks for is available
$23\%$ of the time on commercial $5$G. The obvious fix is predictive --- run a
small action-conditioned world model on both ends and send a frame only once
the prediction has drifted. Two decades of send-on-delta work, and a recent
line of world-model-aided semantic transmission, both do exactly this, and both
decide by \emph{reconstruction fidelity}. We show that fidelity is the wrong
criterion here, because the receiver is not a display but a policy whose
actions determine the next observation. On LIBERO and DROID they are
systematically decoupled: pixel error ranks decision impact at Spearman
$0.31$, and $31.7\%$ of predicted frames are visually wrong yet leave the
policy's action unchanged. A fidelity trigger spends bytes on those frames and
loses tasks to the ones it skips. \sys instead triggers on a
distilled, robot-side estimate of how much a frame would move the policy's
action, and adds the resynchronization protocol that closed-loop predictive
transmission requires and that event-triggered sampling has never had: silence
must be distinguishable from loss. With the same world model and the same task
success rate, \sys sends $59\%$ fewer uplink bytes than fidelity-triggered
transmission and $16\times$ fewer than the deployed JPEG-over-TCP status quo,
and composes with spatial compression for $75\times$ in total.
\end{abstract}

\ccsdesc[500]{Networks~Cyber-physical networks}
\ccsdesc[300]{Networks~Network protocols}
\ccsdesc[300]{Computer systems organization~Robotics}

\keywords{Predictive transmission, event-triggered sampling, world models,
  vision-language-action models, edge offloading, task-oriented communication,
  robot networking, uplink scheduling}

\maketitle

%%%%%%%%%%%%%%%%%%%%%%%%%%%%%%%%%%%%%%%%%%%%%%%%%%%%%%%%%%%%%%%%%%%%%%%%%%%%%%%
\section{Introduction}
\label{sec:intro}

The generalist robot policies that made manipulation broadly capable ---
OpenVLA~\cite{kim2024openvla}, \pizero and \pizerofive~\cite{black2024pi0,
intelligence2025pi05} --- carry $3$--$8$ billion parameters and need tens of
gigabytes of accelerator memory. They do not fit on a robot, and the deployment
that has emerged in response puts the policy on an edge server behind a thin
client: the robot uploads camera frames, the server returns an \emph{action
chunk}, the robot executes it while the next request is in flight. The
reference implementation is explicit about how the observations get there: a
policy server over TCP RPC, with JPEG compression to shrink the
payload~\cite{vlagents2026}. Three camera views at the
resolution the cameras actually produce cost $137$\,kB per upload; at the
$10$\,Hz such a server requests them, that is $11.0$\,Mbps of sustained uplink,
per robot, forever.

Commercial cellular uplink does not supply it. Across the $83$ traces of the
Irish $5$G dataset~\cite{raca20205g} the uplink exceeds $11.0$\,Mbps $23\%$ of
the time and exceeds the $16.4$\,Mbps that a $15$\,Hz observation rate needs
$10.5\%$ of the time (\S\ref{sec:m0}). Downscaling client-side helps and does
not solve it, because the pressure is in the other direction: reactive contact
phases want \emph{higher} observation rates, newer policies want more views,
and a single cell is expected to carry more than one robot. The uplink, not the
GPU and not the downlink, is what limits this deployment.

\paragraph{The obvious fix is old, and it triggers on the wrong quantity.}
Since Miskowicz's send-on-delta
formulation~\cite{miskowicz2006sod}, and in its predictive
variants~\cite{suh2007predictor,nguyen2009stochasticsod} since shortly after,
event-triggered sampling has known that a transmitter which shares a predictor
with its receiver need only speak when the prediction has drifted past a
threshold. Twenty years later the same
idea reappeared with a neural predictor: world-model-aided semantic
transmission runs a generative world model at both ends and \emph{omits
transmission while the prediction remains reliable}, adding a lightweight
feedback module to decide when the current frame must be sent
after all~\cite{wmsemcom2025}. Both lines --- the analytic one
and the generative one --- answer ``should I send this frame?'' with
\emph{how wrong is the receiver's reconstruction}.

That question is correct when the receiver is an estimator with an error
covariance, or a display for a human. It is the wrong question when the
receiver is a policy. A policy does not consume the frame; it consumes the
action the frame implies, and then \emph{acts}, which determines the next
observation. The quantity that matters is not how wrong the predicted frame
looks but whether the true frame would have changed what the robot does. Our
central measurement (\S\ref{sec:m2}) is that these two quantities come apart,
systematically and in both directions. Over $2{,}000$ DROID~\cite{khazatsky2024droid}
episodes and four LIBERO~\cite{liu2023libero} suites, $31.7\%$ of predicted
frames carry large pixel error while leaving the policy's action essentially
unchanged, and $12.4\%$ are pixel-accurate while flipping it; pixel error ranks
decision impact at Spearman $\rho_s = 0.31$. A fidelity trigger spends its
budget on the first group and fails tasks because of the second. It is not
mistuned. It is measuring the wrong thing.

\paragraph{Why the policy is on the server in the first place.}
The premise that the robot can run a world model but not the policy invites an
objection, which we answer by measurement. The gap is not a factor but orders of
magnitude: our world model
costs $3.6$\,GFLOP per control step, against $391$\,GFLOP for the smallest
serious VLA and $4.90$\,TFLOP for OpenVLA-7B, whose $15.2$\,GB of weights do
not fit on the robot at all (Table~\ref{tab:asym}). The split we assume is also
the one the surrounding literature already assumes, independently of us:
latency-tolerant cloud--edge collaborative VLA~\cite{cloudedgevla2026}, remote
policy servers~\cite{vlagents2026}, edge inference
pipelines~\cite{actionflow2025,edgeserving2026,oxygen2026}, and a survey built
on the same deployment constraint~\cite{efmedge2026}. Most importantly the gap
is structural rather than incidental: the useful policy is always chosen larger
than any single robot's budget, while the world model's size is \emph{ours to
choose}. We make this precise with a compute ratio $\rho$ and report where our
mechanism starts and stops paying off (\S\ref{sec:rho}, \S\ref{sec:eval-rho}).

\paragraph{Silence needs a protocol.}
Closing the loop adds a failure mode the open-loop literature does not have. A
silent transmitter is indistinguishable from a lost one --- event-based sampling
work states plainly that it cannot detect dropouts, because the two ends share
no periodic timestamp~\cite{suh2009dropouts}. For a passive receiver that is a
quality-of-reconstruction issue; for a policy it is a divergence hazard, because
the server keeps acting on an imagined observation the robot has already
contradicted, with no alarm. We measure $1.2\%$ uplink loss turning into a
$34.2\%$ episode catastrophe rate under a protocol lacking the distinction
(\S\ref{sec:m4}).

\paragraph{\sys.}
We build the system these measurements imply. A $68$M-parameter
action-conditioned world model runs on both ends, rolled forward on the action
chunks the policy has already emitted, so the two sides stay aligned without
communicating. On the robot, a $2.1$M-parameter \emph{decision-relevance head}
--- distilled offline from the very divergence signal that cannot be computed
online, since computing it would need the $7$B policy the robot cannot host ---
estimates how much the true frame would move the policy's action, and that estimate, not the reconstruction error, is what crosses the
threshold. The threshold itself tracks available uplink, so the mechanism is a
network citizen rather than a fixed-rate consumer. And a two-tier protocol ---
$16$-byte heartbeats carrying sequence numbers and an executed-action hash,
plus a forced full resynchronization --- makes silence unambiguous and bounds
divergence. \sys is orthogonal to spatial compression, which decides how many
bits a frame costs; \sys decides whether the frame is worth sending at all, and
we evaluate the composition.

\paragraph{What we do not claim.}
We did not invent predictive transmission~\cite{miskowicz2006sod,
suh2007predictor}, world-model-aided transmission~\cite{wmsemcom2025}, or
action-conditioned world models~\cite{pilworld2026,hafner2023dreamerv3}: we
take the first two as baselines and reuse the third as a tool. We do not claim
long-horizon prediction --- the mechanism depends on a $5$-step window --- nor
benefits at every compute ratio, and we compose with spatial compression rather
than compete with it.

\paragraph{Contributions.}
(i) A measurement establishing that reconstruction fidelity and decision
relevance are systematically decoupled on embodied observation streams, which
makes fidelity-triggered predictive transmission structurally misallocated in
closed loop and does not depend on our world model being good
(\S\ref{sec:m2}). (ii) A single-variable demonstration that, at the same world model
and the same task success, a decision-relevance trigger sends $59\%$ fewer
uplink bytes than a fidelity trigger (\S\ref{sec:eval-c2}). (iii) A
resynchronization protocol that distinguishes correct silence from loss, with
the divergence bound it enforces (\S\ref{sec:protocol},
\S\ref{sec:eval-proto}). (iv) An explicit regime statement: the compute ratio
$\rho$, its benefit curve across four decades, and both failure boundaries
(\S\ref{sec:eval-rho}).

\begin{figure*}[t]
\centering
\figsch{fig-arch}
\caption{\textbf{\sys's closed-loop predictive transmission path.} The robot
sends a frame only when its distilled head predicts that the true frame would
change the policy's action; heartbeats keep silence distinguishable from loss.
Red marks what \sys adds, grey what it reuses unchanged; \S\ref{sec:design}
walks the numbered path.}
\label{fig:arch}
\end{figure*}

%%%%%%%%%%%%%%%%%%%%%%%%%%%%%%%%%%%%%%%%%%%%%%%%%%%%%%%%%%%%%%%%%%%%%%%%%%%%%%%
\section{Background and Motivation}
\label{sec:motivation}

\subsection{What Is on the Wire, and What the Link Gives}
\label{sec:m0}

\begin{figure}[t]
\centering
\figexp{fig-m0-uplink-cdf}
\caption{\textbf{The deployed configuration does not fit the link.} Uplink CDFs
of three commercial $5$G trace sets against the sustained rate a policy server
needs for three camera views.}
\label{fig:m0}
\end{figure}

A chunked policy maps an observation $o_t$ (two or three camera views plus
proprioception) and a language instruction to a chunk of $H$ future actions,
$H\!\in\![8,50]$~\cite{zhao2023act,chi2023diffusion,black2024pi0}. The robot
executes the chunk at a fixed control rate $f\!=\!15$\,Hz, matching
DROID~\cite{khazatsky2024droid}, while requesting the next one. The asymmetry
between the two directions is severe: the downlink chunk is under $3$\,kB, and
the uplink observation is $137$\,kB for three views at the resolution the
cameras produce. Deployed policy servers upload at that resolution because the
client is a thin RPC wrapper that does not know the policy's preprocessing
pipeline~\cite{vlagents2026}. Downscaling to the policy's $224^2$ input on the
client cuts the payload to $36$\,kB, which is a real improvement and a
one-time constant: it buys back a factor of $3.8$ and then the same argument
applies at $30$\,Hz, or with four cameras, or with four robots on one cell.

Figure~\ref{fig:m0} places the requirement against three commercial trace
sets~\cite{raca20205g,narayanan2020lumos5g,narayanan2021variegated}. The
$11.0$\,Mbps that a $10$\,Hz three-view configuration needs is available $23\%$
of the time on the Irish traces, $45\%$ on Lumos5G --- which is mmWave, and
therefore the optimistic end of what a deployment can hope for --- and $14\%$ on
MN-Wild. At the $30$\,Hz an operator would want for reactive contact phases the
figure is between $1\%$ and $10\%$ everywhere, on every trace set. What we take
from this is not that offloading is impossible --- it plainly works today,
badly --- but that the operating point is set by the link rather than by the
task, and that every byte saved converts into either a higher observation rate
or another robot on the cell. Replaying the status quo over these traces, $41\%$
of uploads queue behind their predecessor and the observation the policy acts on
is on average $310$\,ms old.

\subsection{The Compute Asymmetry Is Real and Structural}
\label{sec:rho}

\begin{table}[t]
\centering
\caption{\textbf{Measured cost of one control step.} Jetson Orin NX ($16$\,GB),
bf16, batch $1$, median of $10^3$ invocations. OpenVLA does not fit in the
board's memory at all.}
\label{tab:asym}
\footnotesize
\setlength{\tabcolsep}{3.4pt}
\begin{tabular}{@{}lcccc@{}}
\toprule
 & \textbf{\sys $W$} & \textbf{SmolVLA} & \textbf{\pizero} &
   \textbf{OpenVLA} \\
 & \textbf{(68\,M)} & \textbf{(450\,M)} & \textbf{(3.3\,B)} &
   \textbf{(7.6\,B)} \\
\midrule
Weights (bf16)        & 0.14\,GB & 0.90\,GB & 6.6\,GB  & 15.2\,GB \\
Peak memory           & 0.21\,GB & 1.4\,GB  & 8.1\,GB  & 16.9\,GB \\
FLOP / control step   & 3.6\,G   & 391\,G   & 2.14\,T  & 4.90\,T \\
\quad ratio to $W$    & $1\times$ & $109\times$ & $594\times$ & $1361\times$ \\
Latency, Orin NX      & 11.7\,ms & 128\,ms  & 1.42\,s$^\dagger$ & OOM \\
Energy / step         & 0.19\,J  & 2.6\,J   & 27\,J$^\dagger$ & --- \\
Latency, A100 server  & 1.9\,ms  & 14\,ms   & 38\,ms   & 62\,ms \\
Meets 66.7\,ms period & \textbf{yes} & no & no & no \\
\bottomrule
\end{tabular}
\\[1mm]
{\scriptsize $^\dagger$ with unified-memory paging; not a deployable
configuration.}
\end{table}

Every predictive-transmission system assumes the transmitter can afford the
predictor. Ours additionally assumes it cannot afford the \emph{consumer}, and
a reviewer is right to ask why. Table~\ref{tab:asym} answers with measurement.
The gap between our world model and the smallest serious VLA is $109\times$ in
per-step arithmetic; against OpenVLA-7B it is $1361\times$, and the more
decisive number is $15.2$\,GB of weights, which excludes essentially every
robot-borne platform. The world model is cheap for a structural reason and not
merely because it is small: its rollout runs in a compact latent space, so
advancing the prediction one step costs $3.4$\,ms and no pixel processing,
while a VLA re-encodes every camera view and runs a language-model backbone on
every query. The robot also never runs the decoder: pixels are reconstructed
only on the server, where the policy needs them.

This is the system model the surrounding literature already adopts, which
matters because rejecting it rejects a subfield rather than a
paper~\cite{cloudedgevla2026,vlagents2026,actionflow2025,edgeserving2026,
efmedge2026}. But the argument we lean on hardest is monotonicity. The frontier
policy is always chosen \emph{larger} than any single robot's budget $B$,
because larger has meant better generalization for three years running; the
world model's size is ours to choose and fits whatever $B$ is. The mechanism is
scale-relative, not scale-absolute:
\begin{equation}
\rho \;=\; \frac{B \cdot 1/f}{C_{\mathrm{VLA}}},
\label{eq:rho}
\end{equation}
the ratio of the robot's per-control-step FLOP budget to the target policy's
per-step demand. On our Orin NX, $B\!\approx\!8.0$\,TFLOP/s sustained and
$C_{\mathrm{VLA}}\!=\!4.90$\,TFLOP give $\rho = 0.109$. We state the regime in
advance and measure both of its edges in \S\ref{sec:eval-rho}: below
$\rho_{\min} = C_W/C_{\mathrm{VLA}} = 7.3\times10^{-4}$ the robot cannot run
the world model within the control period and the mechanism does not apply;
between $\rho_{\min}$ and $1$ is where we claim benefit; at $\rho \ge 1$ the
policy can run locally and the motivation changes character
(\S\ref{sec:discussion}) rather than disappearing.

\subsection{Measurement Setup}
\label{sec:m-setup}

All measurements use public assets. \emph{Offline streams:} $2{,}000$ episodes
sampled across the $86$ tasks of DROID~\cite{khazatsky2024droid} and the four
LIBERO suites~\cite{liu2023libero}, replayed through the policies to collect
observations, action chunks and world-model rollouts; no robot is moved.
\emph{Closed loop:} LIBERO as the primary environment, SimplerEnv in both
Visual Matching and Variant Aggregation settings~\cite{li2024simplerenv}, and
Kinetix~\cite{matthews2025kinetix} as a high-dynamics stress test.
\emph{Policies:} OpenVLA-7B, \pizero, \pizerofive and
SmolVLA~\cite{shukor2025smolvla} at official weights, inference only, never
fine-tuned. \emph{Links:} the Irish $5$G dataset~\cite{raca20205g} as primary
because it publishes uplink throughput, plus
Lumos5G~\cite{narayanan2020lumos5g} and
MN-Wild~\cite{narayanan2021variegated}, replayed with
mahimahi~\cite{netravali2015mahimahi}; controlled loss and delay via
\texttt{tc netem}. \emph{Robot compute:} Jetson Orin NX $16$\,GB.
Inference-time real-time chunking (IT-RTC)~\cite{black2025rtc} is enabled in
every arm, ours and baselines alike, so that no result is attributable to
delay tolerance we did or did not provide.

\subsection{How Far Ahead We Can See}
\label{sec:m1}

\begin{figure}[t]
\centering
\figexp{fig-m1-rollout-error}
\caption{\textbf{Action conditioning roughly doubles the usable prediction
window.} Error is measured in the decision space of \S\ref{sec:trigger},
normalized so that $0.35$ is where the policy's action begins to change
materially (Figure~\ref{fig:m3}b).}
\label{fig:m1}
\end{figure}

Predictive transmission is only interesting if the prediction is usable for
long enough to skip a meaningful number of frames, and our one genuinely
embodiment-specific advantage is that we know the future actions: the policy
has already emitted them, and the downlink has already carried them. B-line
generative transmission cannot use this, because in a video setting there is no
action to condition on --- prediction is driven by text or by visual
extrapolation~\cite{wmsemcom2025,genvidfusion2025}. Figure~\ref{fig:m1}
quantifies what the difference is worth.

Two readings matter. First, the absolute window: with action conditioning,
$63\%$ of LIBERO timesteps stay under the flip threshold for at least $5$ steps
($333$\,ms), comfortably above the $40\%$ we fixed in advance as the minimum
for the mechanism to have anything to work with, and the usable window grows
from $2.9$ steps unconditioned to $5.4$ conditioned; on real DROID streams,
which are teleoperated and therefore jerkier than scripted demonstrations, it
reaches $4.1$. Second, the differential: at $k\!=\!5$ the conditioned
predictor's error is $0.31$ against the unconditioned predictor's $0.58$, a
$46.6\%$ reduction against the $30\%$ we pre-registered as the minimum for
action conditioning to be worth its complexity. This is the entire quantitative
case for doing predictive transmission \emph{in the embodied setting} rather
than importing a video method.

\subsection{Fidelity and Decision Relevance Are Decoupled}
\label{sec:m2}

\begin{figure*}[t]
\centering
\begin{subfigure}{0.48\textwidth}
\centering
\figexp{fig-m2-fidelity-vs-decision}
\end{subfigure}
\hfill
\begin{subfigure}{0.48\textwidth}
\centering
\figexp{fig-m2-distance-ranking}
\end{subfigure}
\caption{\textbf{Fidelity and decision relevance are different quantities.}
(a)~Predicted frames, by pixel reconstruction error against the divergence
$\|\pi(o_t)-\pi(\hat{o}_t)\|$ of the policy's action; a fidelity trigger cuts
vertically, the decision it should make cuts horizontally. (b)~Rank correlation
of six candidate distances with that same ground truth.}
\label{fig:m2}
\end{figure*}

This is the measurement the paper stands on, and it is deliberately independent
of how good our world model is: it asks only which \emph{distance} on predicted
frames ranks their decision impact, which a mediocre predictor can answer. For
every predicted frame from $2{,}000$ DROID episodes and the four LIBERO suites we
compute five candidate distances to the true frame and the ground truth the
online system can never have --- the actual divergence of the policy's action,
$\|\pi(o_t) - \pi(\hat{o}_t)\|$, which costs two \pizero forward passes and is
therefore an offline instrument, not a mechanism. Both thresholds were fixed
before we looked at the data: pixel error above $\rho_s = 0.85$ would have
falsified the paper's premise, our own head below $0.60$ its mechanism.

Figure~\ref{fig:m2} reports the result. Pixel error ranks decision impact at
$\rho_s\!=\!0.31$, and structural and perceptual distances do no better
($0.28$--$0.39$); the world model's own latent distance, which is what a
straightforward adaptation of send-on-delta to high-dimensional observations
would use, reaches $0.52$, and the distilled head of \S\ref{sec:trigger}
reaches $0.74$. The scatter shows why the correlation number
understates the problem. Cut both axes at their operating thresholds and a
fidelity trigger classifies $55.9\%$ of frames the way the policy would;
$31.7\%$ of frames are visually wrong in ways the
policy is indifferent to --- lighting, background texture, the far half of the
scene, the exact pose of an object the gripper has already passed. A fidelity
trigger transmits every one of them. Conversely $12.4\%$ of frames are
reconstructed almost perfectly except in the few pixels that decide the
action: a contact that has or has not occurred, a gripper that has or has not
closed, an object that shifted two millimetres. A fidelity trigger skips every
one of those, and those are the frames whose absence fails the task.

To be precise about what this establishes: not that our trigger is the best
possible one, and nothing about our world model --- a better predictor moves
points toward the origin without rotating the axes. It establishes that the
criterion both the classical and the generative literature use is, in this
regime, measuring a quantity only weakly related to the one that matters, and
that the misallocation is $44.1\%$ of frames, in both directions, rather than a
tail effect. The rest of the paper makes the horizontal cut instead of the
vertical one.

\subsection{The Closed-Loop Cost of an Imagined Observation}
\label{sec:m3}

\begin{figure}[t]
\centering
\begin{subfigure}{0.49\columnwidth}
\centering
\figexp{fig-m3-imagined-steps}
\end{subfigure}
\hfill
\begin{subfigure}{0.49\columnwidth}
\centering
\figexp{fig-m3-decision-error}
\end{subfigure}
\caption{\textbf{The closed-loop price of an imagined observation.}
(a)~Success versus the number of consecutive steps the server spends on
predictions. (b)~The same episodes against the quantity a trigger can actually
observe; the knee at $\tau^\star\!=\!0.35$ is the policy-flip threshold used
throughout.}
\label{fig:m3}
\end{figure}

The previous section says which frames matter. This one says how much slack the
mechanism has, by measuring the closed-loop cost directly rather than inferring
it: we force the server to consume predicted observations for $k$ consecutive
steps and record task success on LIBERO with \pizero, $50$ episodes per point.
Figure~\ref{fig:m3}(a) shows a plateau to about
$4$ steps and a collapse after $8$; panel (b) recasts the same episodes against
the observable quantity, yielding the knee at $\tau^\star = 0.35$ we use
throughout as the level at which the policy starts to behave differently. Two design
consequences follow. The forced-resynchronization interval must be short enough
to keep the tail out of reach, and the trigger threshold must sit far enough
below $\tau^\star$ to absorb the error of the trigger's own estimator --- we
use $\tau_0 = 0.15$ against an estimator whose P99 error is $0.09$.

\subsection{Silence Is Ambiguous, and Dangerously So}
\label{sec:m4}

\begin{figure}[t]
\centering
\figexp{fig-m4-divergence-trace}
\caption{\textbf{What an undetected loss does to a closed loop.} One dropped
resync frame at $t\!=\!6.6$\,s under $1.2\%$ uplink loss, under three
protocols.}
\label{fig:m4}
\end{figure}

The final measurement concerns the failure mode the open-loop literature
identifies and then sets aside. For a passive receiver an undetected loss
degrades a reconstruction; for a policy it is a divergence that compounds rather
than decays, because the server generates actions from an imagined observation
the robot has already contradicted, the robot executes them, and the executed
actions feed both world models.

Figure~\ref{fig:m4} shows one such event under three protocols. We inject
$1.2\%$ uplink loss --- a
mild rate for a congested $5$G uplink --- and drop a single triggered frame.
Without any protocol the divergence never recovers, crosses $\tau^\star$ and
the episode fails $1.8$\,s later; across $200$ episodes at this loss rate,
$34.2\%$ end in a
collision or timeout. Adding periodic forced resynchronization, which is the
natural first fix and the analogue of an I-frame, bounds the excursion but does
nothing about \emph{detection}: the server spends an entire interval acting on
a state it has no reason to distrust, and in $19.7\%$ of episodes that interval
contains a contact event. Sequence numbers, which cost two bytes, turn the same
event into an $84$\,ms excursion. This is why \S\ref{sec:protocol} is a protocol
contribution and not an implementation detail.

\paragraph{Takeaways.} The deployed configuration needs $11$\,Mbps that
commercial $5$G supplies $23\%$ of the time; the robot can afford a predictor
$10^2$--$10^3\times$ cheaper than the policy, structurally so; action
conditioning roughly doubles the usable window; fidelity ranks decision impact
at $\rho_s\!=\!0.31$ and misclassifies $44\%$ of frames in both directions;
imagined observations are free for about $4$ steps and catastrophic after $10$;
and unprotected silence turns $1.2\%$ loss into $34.2\%$ episode failures.

%%%%%%%%%%%%%%%%%%%%%%%%%%%%%%%%%%%%%%%%%%%%%%%%%%%%%%%%%%%%%%%%%%%%%%%%%%%%%%%
\section{Design}
\label{sec:design}

\sys has four parts: a shared action-conditioned world model
(\S\ref{sec:wm}), a decision-relevance trigger that replaces the fidelity test
(\S\ref{sec:trigger}), a bandwidth-adaptive threshold (\S\ref{sec:adaptive}),
and the resynchronization protocol (\S\ref{sec:protocol}).
Figure~\ref{fig:arch} shows how they compose, and its numbered path is one
control step. \cnum{1}~The robot encodes the true observation. \cnum{2}~Both
ends roll the identical world model forward on the action chunks the policy has
already emitted, so they stay aligned while the uplink is silent. \cnum{3}~A
distilled head estimates how much the discrepancy between the two would move
the policy's action, and \cnum{4}~only that estimate --- never a reconstruction
error --- crosses the threshold. \cnum{5}~When no frame arrives the server
decodes $\hat{o}_t$ for the policy, and \cnum{6}~the returned chunk drives the
executor and, through it, both world models. The robot never runs the decoder
and never runs the policy. The world model is not a contribution:
action-conditioned predictors that generate policy-consumable observations from
a rollout of committed actions already exist and match a policy's input
format~\cite{pilworld2026,hafner2023dreamerv3,wu2024ivideogpt}. We adopt one,
shrink it, and claim nothing about it; our contributions are the trigger and the
protocol, both indifferent to which predictor sits underneath.

\subsection{The Shared Predictor}
\label{sec:wm}

Both ends run identical weights $W$: an encoder $E$ mapping three camera views
to a $96$-token latent $z_t$, a chunk-wise latent dynamics model
$\hat{z}_{t} = W(\hat{z}_{t-1}, a_{t-1})$, and --- on the server only --- a
decoder $D$ that renders $\hat{o}_t$ in the format the policy expects. Total
$68$M parameters, $0.14$\,GB of weights, $11.7$\,ms per control step on the
robot (Table~\ref{tab:asym}). We chose the size to serve the argument of
\S\ref{sec:rho} rather than to maximize prediction quality: a smaller model
makes the compute-asymmetry claim stronger, and neither the decoupling result
nor the protocol depends on prediction quality. A large video diffusion model
would predict better, fail to run on the robot, and invite the objection that
the robot might as well run the policy.

The two ends agree because they share weights, an initial state and an action
sequence: the actions travel down the link
and the robot knows what it executed. Consistency in floating point does not
follow from consistency in algorithm, and this bit us: the same graph on an Orin
NX and on an A100 drifts apart by $4\times10^{-4}$ per step in fp16. We quantize
the recurrent latent to $8$ bits after every step, which makes the rollout
exactly reproducible across platforms for $0.4$\pp of success.

\subsection{Triggering on Decision Relevance}
\label{sec:trigger}

\paragraph{The quantity we want, and why we cannot have it.}
The ideal trigger fires when the true observation would change the policy's
action, i.e.\ on $\delta_t = \|\pi(o_t) - \pi(\hat{o}_t)\|$. Computing it
requires two forward passes of the policy on the robot, which is exactly the
thing the robot cannot do --- this is not an efficiency concern but a
definitional one, since a robot that could evaluate $\pi$ twice per step would
not be offloading. So $\delta_t$ is available offline, in quantity, and never
online.

\paragraph{The distilled head.} We therefore train a small head $\Phi$ to
predict it. Given the encoded true observation $z_t$ and the predicted latent
$\hat{z}_t$, both of which the robot already has,
\begin{equation}
\hat{\delta}_t \;=\; \Phi\big(z_t,\, \hat{z}_t,\, z_t - \hat{z}_t\big),
\label{eq:head}
\end{equation}
a $2.1$M-parameter MLP over pooled latents, trained by regression on
$\delta_t$ computed offline over LIBERO and DROID with the policy in the loop
($1.4$M frame pairs, $9$ GPU-hours, once). At runtime it costs $1.1$\,ms and
never touches the policy. A concurrent submission applies the same methodology
--- online-unavailable ground truth, obtainable proxy, offline calibration ---
to a different quantity in the same deployment~\cite{policyvoi2027}; the two are
orthogonal, latency against bandwidth.

An obvious alternative is to compare in the policy's own visual-encoder latent
space, which is decision-relevant by construction. We evaluate it and do not use
it: a VLA's vision tower costs $214$\,GFLOP per step, $59\times$ the rest of our
pipeline, which would destroy the argument of \S\ref{sec:rho}. The head is a
$200\times$ cheaper approximation, and the gap is $1.6$\pp.

\paragraph{The rule.} The robot transmits when
\begin{equation}
\hat{\delta}_t > \tau(\hat{B}_t)
\quad\text{or}\quad
k_t \ge T_{\mathrm{resync}}
\quad\text{or}\quad
\text{NACK received},
\label{eq:trigger}
\end{equation}
where $k_t$ counts steps since the last transmitted frame. The three clauses
are the mechanism, the safety net, and the protocol respectively, and they are
deliberately not merged: the first is an estimate and may be wrong, the second
bounds how wrong it can be for how long, the third handles the link.

\subsection{Making the Threshold a Network Citizen}
\label{sec:adaptive}

A fixed $\tau$ would make \sys a source that ignores the link, which is the
mistake the status quo already makes. We set
\begin{equation}
\tau(\hat{B}_t) = \mathrm{clip}\!\left(\tau_0
  \left(\frac{B_{\mathrm{ref}}}{\hat{B}_t}\right)^{\!\alpha},
  \tau_{\min}, \tau_{\max}\right),
\label{eq:tau}
\end{equation}
with $\tau_0 = 0.15$, $\alpha = 0.5$, $[\tau_{\min},\tau_{\max}] =
[0.08, 0.30]$ and $\hat{B}_t$ an EWMA of delivered uplink throughput. When the
link is scarce the threshold relaxes and the robot speaks only about frames
that would definitely change the action; when the link is generous it tightens
and the system approaches continuous streaming. The clip at $\tau_{\max}$ is
what keeps the mechanism honest: the threshold is never allowed above $0.30$,
which is below the policy-flip level $\tau^\star = 0.35$ of
Figure~\ref{fig:m3}(b), so degradation under congestion is bounded by
construction rather than by tuning.

\subsection{The Resynchronization Protocol}
\label{sec:protocol}

\begin{figure}[t]
\centering
\figsch{fig-proto}
\caption{\textbf{The protocol makes silence unambiguous.} Heartbeats carry a
sequence number and a rolling hash of the actions the robot actually executed,
so the server can separate correct silence from loss and demand recovery.}
\label{fig:proto}
\end{figure}

Figure~\ref{fig:proto} gives the state machine. Three mechanisms carry it.
\emph{Sequence numbers} on every transmitted frame let the server detect a gap
and issue a NACK, which is the distinction \S\ref{sec:m4} shows is missing
from published event-triggered designs. \emph{Heartbeats} every $T_h = 8$
control steps carry the current sequence number, the robot's own
$\hat{\delta}_t$, and a rolling hash of the last eight executed actions; the
hash catches the second divergence source, in which a lost or late
\emph{downlink} chunk makes the robot fall back to its safety behaviour so that
the actions it executed are not the actions the server assumed. \emph{Forced
resynchronization} every $T_{\mathrm{resync}} = 60$ steps ($4$\,s) bounds drift
in the case the trigger itself is wrong, which the head's finite accuracy
guarantees will happen; it is the I-frame of this scheme, and its cost is a
$1.7\%$ floor on the send rate.

The resulting divergence is empirically bounded by
\begin{equation}
d_t \;\le\; \tau(\hat{B}_t) \;+\; \varepsilon_\Phi \;+\;
  \beta \cdot \min\!\big(k_t,\, T_{\mathrm{resync}}\big)\,
  \mathbf{1}[\Phi \text{ misses}],
\label{eq:bound}
\end{equation}
where $\varepsilon_\Phi$ is the head's error (P99 $= 0.09$) and
$\beta = 0.07$/step the measured open-loop drift rate from
Figure~\ref{fig:m1}. The head misses on $1.3\%$ of steps, so the third term is
rare; the design rule is to pick $T_{\mathrm{resync}}$ so that the bound stays
under $\tau^\star$ at the offline-measured miss rate. \S\ref{sec:eval-proto}
reports P99.9 divergence of $0.21$ against a predicted bound of $0.24$.

\subsection{Composition, and What Breaks}
\label{sec:compose}

\sys occupies the temporal axis --- whether a frame is sent --- and is
orthogonal to the spatial axis, which decides what a sent frame costs. We run
IT-RTC~\cite{black2025rtc} underneath everything as a latency substrate, and we
compose with a learned VLA-specific frame codec~\cite{licvla2026} on triggered
frames; \S\ref{sec:eval-main} evaluates both compositions and neither
interferes.

Three failure modes are handled explicitly. On \emph{cold start} the first two
frames of every episode are transmitted unconditionally. If the \emph{head is
unavailable or its input stale}, the trigger degrades to
$\hat{\delta}_t = \infty$ and the system falls back to continuous transmission
--- to the status quo, not to something worse. And if the \emph{uplink cannot
carry even the forced resynchronizations}, measured at sustained uplink below
$0.6$\,Mbps, the robot stops rather than act on an unbounded-divergence plan;
this is the convention production teleoperation stacks use and it is enabled in
every arm of our evaluation, baselines included.

%%%%%%%%%%%%%%%%%%%%%%%%%%%%%%%%%%%%%%%%%%%%%%%%%%%%%%%%%%%%%%%%%%%%%%%%%%%%%%%
\section{Implementation}
\label{sec:impl}

\sys is $5{,}200$ lines of Python and $1{,}100$ lines of C++, and reuses four
public codebases rather than reimplementing them.

\paragraph{World model and policies.} The predictor follows the chunk-wise
action-conditioned architecture of policy-in-the-loop world
models~\cite{pilworld2026}, reduced to $68$M parameters ($6$ latent-transformer
blocks, width $384$, $96$ tokens per step) and trained on the $2{,}000$ LIBERO
demonstrations plus a $40$k-episode DROID subset for $19$ hours on eight
RTX~3090s. Policies are served through LeRobot~\cite{cadene2024lerobot}, whose
\texttt{policies.rtc} module provides the reference IT-RTC implementation for
\pizero, \pizerofive and SmolVLA; OpenVLA-7B is wrapped behind the same
interface. The server is organized after the unified policy-server design of
VLAgents~\cite{vlagents2026}, one process per GPU with a shared request queue,
which is also our status-quo baseline.

\paragraph{Robot client.} The encoder, rollout, head and trigger run in the
control thread on a Jetson Orin NX as a single TensorRT engine with the
$8$-bit latent quantization of \S\ref{sec:wm} applied between steps; this is
the C++ part, because it is the only code with a real-time deadline. The client
holds no policy weights.

\paragraph{Transport.} The uplink is a userspace UDP path with a $16$-byte
control header carrying sequence number, action hash and $\hat{\delta}_t$, and
a JPEG or learned-codec payload for triggered frames. NACKs and forced-resync
requests travel on the downlink alongside action chunks. Trace replay uses
mahimahi~\cite{netravali2015mahimahi} with per-direction link shells driven by
the Irish, Lumos5G and MN-Wild uplink series; loss and delay injection uses
\texttt{tc netem}. We deliberately do not run over TCP, unlike the status
quo~\cite{vlagents2026}: head-of-line blocking on a link this scarce turns one
lost frame into a multi-hundred-millisecond stall, and our protocol already
provides the loss detection that TCP would be supplying.

\paragraph{Compute-budget emulation.} The $\rho$ sweep needs robot budgets we
do not own hardware for. We emulate them by pinning the client to a fraction of
the Orin's SMs via MPS and capping the clock, calibrating each setting against
measured FLOP/s so that the reported $\rho$ is achieved rather than nominal.
The $\rho = 0.109$ point is native hardware and anchors the curve.

%%%%%%%%%%%%%%%%%%%%%%%%%%%%%%%%%%%%%%%%%%%%%%%%%%%%%%%%%%%%%%%%%%%%%%%%%%%%%%%
\section{Evaluation}
\label{sec:eval}

We ask five questions: does a decision-relevance trigger beat a fidelity trigger
when nothing else differs (\S\ref{sec:eval-c2}); where does \sys sit against the
full range of alternatives (\S\ref{sec:eval-main}); over what compute ratios
does the mechanism pay (\S\ref{sec:eval-rho}); does the protocol do what
\S\ref{sec:protocol} claims (\S\ref{sec:eval-proto}); and what does it cost, in
simulation and on hardware (\S\ref{sec:eval-ablation}--\S\ref{sec:eval-real})?

\paragraph{Setup and baselines.} Environments, policies, traces and hardware
are as in \S\ref{sec:m-setup}; every configuration runs at least $50$ episodes,
error bars are $95\%$ bootstrap CIs, and IT-RTC is enabled everywhere. An
episode is $50$\,s, $750$ control steps, and a transmitted frame is $137$\,kB
unless a codec says otherwise. The baselines fall into three groups, and
the grouping is the argument. \emph{Status quo and conventional coding:} (1) the
deployed policy server, fixed $10$\,Hz JPEG over TCP~\cite{vlagents2026}; (2)
fixed $10$\,Hz H.265 with rate control~\cite{sullivan2012hevc}; (3)
pixel-difference change detection. \emph{Classical predictive transmission:} (4)
send-on-delta with a Taylor predictor on the proprioceptive
state~\cite{miskowicz2006sod,suh2007predictor}; (5) the same rule lifted to the
image latent. \emph{Modern generative and spatial competitors:} (6) a
world-model transmitter with a \emph{fidelity} trigger --- our reproduction
of~\cite{wmsemcom2025} on \emph{the identical world model as ours}; (7) learned
image compression for VLA~\cite{licvla2026}; (8) \sys; (9) \sys composed with
(7).

\subsection{Fidelity vs.\ Decision: The Single-Variable Result}
\label{sec:eval-c2}

\begin{figure}[t]
\centering
\figexp{fig-c2-byte-reduction}
\caption{\textbf{The one comparison that isolates the paper's mechanism.}
Identical world model, policy, link and protocol; each arm's threshold is tuned
so that the two match on task success within $0.7$\pp.}
\label{fig:c2}
\end{figure}

We put this first because it is the paper's only clean evidence for the
mechanism; every other comparison confounds it with something else.
Figure~\ref{fig:c2} holds the world model, the policy, the
link, the protocol and the success rate fixed, and varies only the quantity
that crosses the threshold. Averaged over the four LIBERO suites, triggering on
estimated decision impact rather than on reconstruction error sends $59.2\%$
fewer uplink bytes at the same task success.

The number is the direct consequence of Figure~\ref{fig:m2}: the fidelity
trigger transmits the $31.7\%$ of frames that are visually wrong and
behaviourally irrelevant, and \sys does not. Worth noting is what \emph{did
not} happen --- success did not fall. A trigger that merely sent less would have
paid for it; \sys does not, because it spends part of the saving on the $12.4\%$
of frames a fidelity trigger skips, so its send rate ($4.1\%$ of control steps)
is a differently-shaped set rather than a scaled-down version of the fidelity
trigger's ($10.0\%$). Kinetix is the honest boundary case: its dynamics shorten
the usable window to $2.6$ steps, both arms transmit far more often, and the
advantage falls to $46.7\%$. We expect any domain with sub-$200$\,ms scene
dynamics to behave this way.

\subsection{The Full Frontier}
\label{sec:eval-main}

\begin{table*}[t]
\centering
\caption{\textbf{End-to-end results.} Send rate is the \% of control steps at
which a frame is transmitted; divergence is measured in the decision space of
\S\ref{sec:trigger}, where $0.35$ flips the action.}
\label{tab:main}
\scriptsize
\begin{tabular}{@{}llccccccc@{}}
\toprule
& \textbf{Transmission policy} & \textbf{LIBERO} & \textbf{SimplerEnv} &
  \textbf{Kinetix} & \textbf{Uplink} & \textbf{Send rate} &
  \textbf{P99.9 div.} & \textbf{vs.\ status quo} \\
\midrule
& Unlimited-uplink oracle (upper bound) & 79.4 & 70.6 & 68.2 & --- & 100\% &
  --- & --- \\
\midrule
\multirow{3}{*}{\rotatebox{90}{\tiny status quo}}
& (1) Policy server, 10\,Hz JPEG/TCP~\cite{vlagents2026} & 71.2 & 62.8 & 55.1 &
  68.4 & 66.7\% & --- & $1\times$ \\
& (2) Fixed 10\,Hz H.265, rate-controlled & 74.6 & 65.9 & 58.7 & 11.9 & 66.7\% &
  --- & $5.7\times$ \\
& (3) Pixel-difference change detection & 72.9 & 64.1 & 57.0 & 9.4 & 9.1\% &
  0.61 & $7.3\times$ \\
\midrule
\multirow{2}{*}{\rotatebox{90}{\tiny classic}}
& (4) Send-on-delta + Taylor, proprioceptive~\cite{suh2007predictor} & 68.7 &
  60.3 & 52.4 & 14.7 & 14.3\% & 0.74 & $4.7\times$ \\
& (5) Send-on-delta lifted to image latent & 73.5 & 65.0 & 58.1 & 12.3 &
  12.0\% & 0.44 & $5.6\times$ \\
\midrule
\multirow{4}{*}{\rotatebox{90}{\tiny modern}}
& (6) World model + fidelity trigger~\cite{wmsemcom2025} & 78.2 & 69.1 & 63.4 &
  10.3 & 10.0\% & 0.29 & $6.6\times$ \\
& (7) Learned image compression for VLA~\cite{licvla2026} & 75.8 & 67.2 & 60.8 &
  7.7 & 66.7\% & --- & $8.9\times$ \\
& \textbf{(8) \sys} & \textbf{78.9} & \textbf{69.8} & \textbf{64.6} &
  \textbf{4.2} & \textbf{4.1\%} & \textbf{0.21} & $\mathbf{16.3\times}$ \\
& \textbf{(9) \sys $+$ (7)} & 78.4 & 69.4 & 64.0 & \textbf{0.9} & 4.1\% & 0.21 &
  $\mathbf{75.2\times}$ \\
\bottomrule
\end{tabular}
\end{table*}

\begin{figure}[t]
\centering
\figexp{fig-pareto}
\caption{\textbf{The byte--success frontier.} \sys dominates every baseline at
every operating point; each curve is a sweep of that arm's own threshold.}
\label{fig:pareto}
\end{figure}

Table~\ref{tab:main} and Figure~\ref{fig:pareto} place that mechanism in
context. Both report \pizero with IT-RTC over the Irish $5$G traces; uplink is
MB per episode and the LIBERO column averages the four suites. Four
observations. First,
the classical baselines behave exactly as
their assumptions predict: send-on-delta on the proprioceptive state (4) is the
worst arm in the table, because a threshold on joint positions is blind to
everything the cameras exist to see, and lifting the same rule to the image
latent (5) recovers most of the loss but inherits the fidelity criterion and with
it the misallocation of \S\ref{sec:m2}. Second, conventional video coding (2) is
stronger than one might expect --- H.265 exploits the same temporal redundancy
with no model at all --- and remains $2.8\times$ more expensive than \sys
because it still pays for every frame. Third, the spatial competitor (7) is
genuinely good and genuinely orthogonal, and the combination (9) reaches
$0.9$\,MB per episode, $75\times$ below the status quo, at a $0.5$\pp success
cost we attribute to codec artifacts on the frames that matter most.

Fourth, and least expected: \sys \emph{improves} on the status quo's success
rate by $7.7$\pp while sending $16\times$ less. The status quo is not merely
wasteful but self-harming --- at $11$\,Mbps demanded from a link that supplies
it $23\%$ of the time, $41\%$ of its uploads arrive stale (\S\ref{sec:m0}).
Sending less, when what you send is chosen well, makes it arrive on time.

The frontier adds what the table cannot show: across the whole sweep the two
world-model arms stay separated by a roughly constant factor of $2.4$ in bytes,
the signature of a criterion difference rather than of a tuning difference --- no
threshold setting moves the fidelity arm onto our curve. We operate at the
$4.2$\,MB knee; pushing $\tau$ below $0.15$ buys $0.3$\pp for $55\%$ more
bytes.

\subsection{Where the Mechanism Applies}
\label{sec:eval-rho}

\begin{figure}[t]
\centering
\figexp{fig-rho-compute-ratio}
\caption{\textbf{Benefit versus compute ratio, with both failure boundaries.}
Shading marks the region in which the world model cannot finish inside the
control period and \sys degrades to fixed-rate transmission.}
\label{fig:rho}
\end{figure}

Figure~\ref{fig:rho} answers the question this system model always provokes.
The benefit is not a knife-edge property of one hardware configuration: it stays
above $50\%$ across two and a half decades of $\rho$ and saturates rather than
peaks, because past a point the world model is good enough and more compute buys
prediction quality the decision trigger does not need. The rise is steepest up to
$\rho \approx 10^{-2}$, the point at which a model large enough for a $5$-step
window fits inside the budget. The lower boundary is
sharp: at $\rho < \rho_{\min} = 7.3\times10^{-4}$ the predictor misses the
$66.7$\,ms control deadline
and the right engineering answer is a fixed-rate encoder, not this paper. The
upper boundary is soft and is not where the value goes to zero --- at $\rho = 1$
the byte reduction is still $63\%$, and only the reason to want it changes
(\S\ref{sec:discussion}). What the sweep says, in one sentence, is that the
mechanism is scale-relative: what matters is that the policy is larger than the
robot's budget, not how large either one is.

\subsection{Does the Protocol Earn Its Complexity?}
\label{sec:eval-proto}

We sweep uplink loss from $0$ to $5\%$ with \texttt{tc netem} at a fixed
trigger. Without sequence numbers or forced resynchronization --- the
configuration published event-triggered designs describe --- the episode
catastrophe rate (collision or timeout) rises from $1.8\%$ at zero loss to
$34.2\%$ at $1.2\%$ loss and $71.5\%$ at $5\%$. Forced resynchronization alone
cuts these to $9.4\%$ and $22.1\%$: it bounds the damage but cannot detect it,
so the server acts on a diverged state for up to a full interval. The full
protocol reaches $2.1\%$ and $4.6\%$ with a median detect-to-recover time of
$84$\,ms, for $240$\,B/s of heartbeat traffic --- $0.02\%$ of what it
protects.

The divergence bound of Equation~\ref{eq:bound} also holds: over $2{,}400$
episodes the P99.9 decision-space divergence is $0.21$ against a predicted
bound of $0.24$ and a policy-flip threshold of $0.35$, and the maximum observed
value, $0.33$, occurred in the one configuration where the head's miss rate
exceeded its offline estimate. Sensitivity to $T_{\mathrm{resync}}$ is mild
between $30$ and $120$ steps ($\pm 0.9$\pp success, $\pm 1.4\%$ bytes) and
degrades sharply past $180$ steps, where Equation~\ref{eq:bound} predicts the
bound crossing $\tau^\star$ and success duly falls by $6.1$\pp.

\subsection{Ablations}
\label{sec:eval-ablation}

\begin{figure}[t]
\centering
\figexp{fig-ablation}
\caption{\textbf{Ablation} at $1.2\%$ uplink loss, so that the protocol
components are exercised. Purple marks the two protocol rows, which lose success
\emph{without} saving bytes.}
\label{fig:ablation}
\end{figure}

Figure~\ref{fig:ablation} separates the components, and the ordering is
informative. The largest single loss comes from removing sequence numbers
($-17.7$\pp), which is not a performance knob but the difference between a
protocol that can detect divergence and one that cannot; removing action
conditioning costs $7.6$\pp and $64\%$ more bytes, confirming
Figure~\ref{fig:m1}. Replacing the distilled head with the world model's own
latent distance --- the most natural way to build this system without our
measurement --- costs $2.5$\pp and $69\%$ more bytes, and replacing it with a
fidelity criterion costs $145\%$ more bytes, the same $59\%$ reduction seen
from the other direction. A static $\tau$ costs $4.3$\pp and $38\%$ more bytes.

Two rows must be read with their uplink volumes, which the figure does not draw.
The protocol ablations send $4.0$ and $3.4$\,MB against the full system's $4.2$,
so their $17.7$ and $10.0$\pp losses buy nothing --- a correctness bug rather
than a trade-off, and the reason we treat the protocol as part of the mechanism.
The oracle row, which substitutes the ground-truth divergence of
\S\ref{sec:m2} for the head, bounds what a better estimator could add at
$0.4$\pp and $10\%$ of bytes: the head is close enough to its own ceiling that
the headroom is elsewhere.

Three further sweeps. \emph{World-model size:} at $18$M the prediction window
shrinks to $3.1$ steps and success falls to $76.1\%$; at $210$M it reaches
$79.2\%$ but misses the control period on the Orin --- precisely the trade
$\rho$ describes. \emph{Latent quantization:} the $8$-bit state that buys exact
cross-platform reproducibility costs $0.4$\pp, and without it $3.1\%$ of
episodes accumulate silent numerical divergence. \emph{Head transfer:} a head
distilled on \pizero and applied unchanged to OpenVLA-7B costs $2.3$\pp and
$8\%$ more bytes, and one trained jointly on three policies recovers all but
$0.6$\pp: the calibration is policy-sensitive, not policy-specific.

\subsection{Cost on the Robot}
\label{sec:eval-overhead}

\begin{table}[t]
\centering
\caption{\textbf{Robot-side cost per control step}, Orin NX, median of $10^5$
invocations; the $15$\,Hz control period is $66.7$\,ms.}
\label{tab:overhead}
\footnotesize
\setlength{\tabcolsep}{4pt}
\begin{tabular}{@{}lccc@{}}
\toprule
\textbf{Component} & \textbf{Latency} & \textbf{Memory} & \textbf{Energy} \\
\midrule
Encoder $E$, three views       & 6.9\,ms & 0.11\,GB & 0.11\,J \\
Latent rollout $W$, one step   & 3.4\,ms & 0.06\,GB & 0.05\,J \\
Decision head $\Phi$ + trigger & 1.1\,ms & 0.02\,GB & 0.02\,J \\
Protocol (seq, hash, framing)  & 0.3\,ms & $<$0.01\,GB & 0.01\,J \\
\midrule
\textbf{Total}                 & \textbf{11.7\,ms} & \textbf{0.21\,GB} &
  \textbf{0.19\,J} \\
\quad as \% of control period  & 17.5\% & --- & --- \\
\bottomrule
\end{tabular}
\end{table}

Table~\ref{tab:overhead} reports the price of running a predictor on the robot,
in full, because the honest accounting is less flattering than the byte numbers.
\sys costs $17.5\%$ of the control period and $2.85$\,W of sustained compute. It
saves radio energy --- the status quo's $10$ uploads per second cost $5.2$\,W of
$5$G transmission against $0.32$\,W for ours --- so the net is $5.5$\,W down to
$3.2$\,W, a $1.7\times$ reduction rather than the $16\times$ the byte numbers
suggest. Evaluating the trigger at $7.5$\,Hz halves the compute to $1.42$\,W for
$0.4$\pp of success, bringing the net to $2.6\times$. \sys trades robot compute
for uplink capacity, and on a shared cell the uplink is the scarce resource, not
the robot's watts.

\subsection{Physical-Robot Evaluation}
\label{sec:eval-real}

\begin{figure}[t]
\centering
\figexp{fig-real-franka}
\caption{\textbf{Franka Panda, $160$ trials.} Two tasks against five shaped
uplink capacities; bars are $95\%$ CIs.}
\label{fig:real}
\end{figure}

Simulated contact dynamics are only as trustworthy as the simulator, so we ran a
controlled hardware experiment on a Franka Panda --- the DROID platform, chosen
so that the offline analysis and the hardware share an embodiment --- with two
tasks (peg insertion and cluttered pick-and-place), five shaped uplink capacities
and $160$ trials. This is a variable sweep,
not a data collection: nothing is trained and no dataset is built. The result we
care about in Figure~\ref{fig:real} is not the peak but the shape. The two arms
converge above $8$\,Mbps, as they must; below $4$\,Mbps the status quo degrades
steeply while \sys stays nearly flat, because its send rate is set by what the
scene does rather than by a clock. At $2$\,Mbps --- a congested commercial
uplink --- the status quo succeeds on $35\%$ of trials and \sys on $75\%$,
sending $0.44$\,MB per episode against $5.8$\,MB. That measured byte ratio,
$13.2\times$, is within $20\%$ of simulation's $16.3\times$.

%%%%%%%%%%%%%%%%%%%%%%%%%%%%%%%%%%%%%%%%%%%%%%%%%%%%%%%%%%%%%%%%%%%%%%%%%%%%%%%
\section{Related Work}
\label{sec:related}

\paragraph{Send-on-delta and event-triggered sampling.} We did not invent this
mechanism class. Lebesgue sampling~\cite{astrom2002lebesgue,arzen1999pid} and
Miskowicz's send-on-delta~\cite{miskowicz2006sod} established transmitting on
signal deviation rather than on a clock; predictive variants transmit when the
signal departs from an extrapolation the receiver can reproduce, with truncated
Taylor~\cite{suh2007predictor} or linear~\cite{nguyen2009stochasticsod}
predictors, and derive thresholds from error
bounds~\cite{miskowicz2007kldiv}; controlled-communication policies solve for an
optimal transmission policy against a local copy of the remote
estimate~\cite{molin2013optimal,mamduhi2017error,
heemels2012introduction,tabuada2007eventtriggered}. Every element of that
setting differs from ours in the same direction: the payload is a
low-dimensional state, not multi-camera video; the predictor is analytic, not
learned; the receiver is an estimator with an error covariance, not a black-box
$7$B policy; and the trigger comes from an estimation bound rather than from a
decision. We take this line as baseline group (4)--(5) rather than as
background, and inherit its open problem: that event-based sampling cannot
detect dropouts for want of a shared periodic timestamp~\cite{suh2009dropouts}
is what \S\ref{sec:protocol} addresses.

\paragraph{World models and generative semantic communication.} The nearest
neighbour is world-model-aided semantic transmission, which runs a generative
world model at both ends, omits transmission while prediction stays reliable,
and adds a feedback module to decide when the current frame must be
sent~\cite{wmsemcom2025}, and a broader literature transmits sketches,
keyframes or tokens and reconstructs generatively~\cite{gosc2025video,
genvidfusion2025,resividtok2026,agentsemcom2025,genvidcomp2025,gunduz2023beyond,
kountouris2021semantics,bourtsoulatze2019deepjscc}. What
unites them is a \emph{passive} receiver: the reconstruction is rendered, scored
by PSNR or perceptual or semantic consistency, and does not act. Ours acts, and
its actions determine the next observation, which is what makes reconstruction
fidelity the wrong trigger (\S\ref{sec:m2}) rather than an imperfect one. We
reproduce \cite{wmsemcom2025}'s design as baseline (6), with our world model,
so that the comparison isolates the criterion.

The exception is a world-model digital twin that selects semantic tokens by
their marginal contribution to expected long-term return, normalized per
bit~\cite{cdtsemcom2026}: the closest anyone comes to our criterion, and
explicit that a one-shot fidelity metric is the wrong one for a closed loop. It
differs in what it estimates and in what it decides --- a learned long-horizon
return, from an actor--critic it trains inside the twin, against our one-step
divergence in the action space of a policy we never retrain; and \emph{which
tokens} a transmission carries, against \emph{whether to transmit at all}, the
decision that makes silence ambiguous and forces \S\ref{sec:protocol}.

\paragraph{Predictive display and model-mediated teleoperation.} Predicting the
remote scene to hide delay is the oldest line here, from the phantom
robot~\cite{bejczy1990phantom} and passivity-based teleoperation~\cite{
anderson1989bilateral,hannaford2002passivity} through environment-adapted
controllers~\cite{passenberg2010survey,brudnak2016predictive} to recent
zero-shot benchmarks of video models as predictive
displays~\cite{genpreddisplay2026}. The prediction there is for a human, the
purpose is to compensate perceived latency, and the evaluation is operator
performance and cognitive load. We predict for a model, to avoid sending bytes,
which makes our criterion automatable and calibratable in a way theirs cannot
be.

\paragraph{Observation compression for VLA.} A fast-growing line compresses
what an offloaded policy consumes. Learned image compression tuned for VLA
notes that general-purpose codecs optimize perceptual fidelity rather than
downstream control and that visual importance is highly non-uniform across
cameras and regions~\cite{licvla2026}; token-level methods distil each frame to
a single latent token~\cite{onewmvla2026}, cache or prune visual
tokens~\cite{compressorvla2025,vlacache2025}, or compress task-oriented features
for device--edge and vehicle-to-everything collaboration~\cite{tofc2025,
dinolink2026}. All
allocate bits \emph{within} a frame; none decides whether the frame is worth
sending. The axes are orthogonal and Table~\ref{tab:main} row (9) shows they
multiply. The same holds for video-analytics systems that filter on-camera or
let the server drive region
selection~\cite{li2020reducto,du2020dds,zhang2018awstream,chen2015glimpse}: their
consumer scores per-frame accuracy, not a closed loop in which its next input is
a consequence of its own output.

\paragraph{Action-conditioned world models.} Policy-in-the-loop world models
generate futures chunk by chunk from a policy's own action rollout and match the
policy's input format~\cite{pilworld2026}, building on a line of learned
simulators~\cite{ha2018worldmodels,hafner2023dreamerv3,bruce2024genie,
wu2024ivideogpt}. We use one as a component: our results are stated relative to
a fixed predictor and would move with a better one without changing sign, the
property \S\ref{sec:m2} was designed to guarantee.

\paragraph{Edge serving and latency tolerance for VLA.} Splitting inference
between device and edge is long established~\cite{kang2017neurosurgeon}; recent
systems specialize it for policies with latency-tolerant cloud--edge
collaboration~\cite{cloudedgevla2026}, policy
servers~\cite{vlagents2026}, KV-cache management~\cite{oxygen2026}, pipelined
edge inference~\cite{actionflow2025} and deadline-aware
scheduling~\cite{edgeserving2026}, all under the deployment constraint a recent
survey takes as given~\cite{efmedge2026}. A parallel line makes the policy
tolerate delay rather than reduce it~\cite{black2025rtc,vlash2025,
asyncbench2026}; we run IT-RTC underneath every arm because it is a substrate,
not a competitor. Closest in method is a concurrent submission that reads a
value-of-information signal off a policy's own action chunks to decide
\emph{when} an inference is worth requesting~\cite{policyvoi2027}: same
calibration methodology, latency axis instead of bandwidth.

%%%%%%%%%%%%%%%%%%%%%%%%%%%%%%%%%%%%%%%%%%%%%%%%%%%%%%%%%%%%%%%%%%%%%%%%%%%%%%%
\section{Discussion and Limitations}
\label{sec:discussion}

\paragraph{The upper bound is the world model.} Everything \sys saves is
bounded by how long the shared predictor stays useful, and
Figure~\ref{fig:c2}'s Kinetix column is the honest version of this: when scene
dynamics outrun the prediction window, both triggers converge toward continuous
transmission and our advantage shrinks from $59\%$ to $47\%$. We expect the
mechanism to suit fast locomotion and scenes with independently moving agents
poorly --- precisely the settings where the action conditioning of
Figure~\ref{fig:m1} stops paying.

\paragraph{At $\rho \ge 1$ the mechanism survives; the motivation changes.}
Figure~\ref{fig:rho} shows the byte reduction holding at $63\%$ when the robot
could run the policy locally, and it is fair to ask why anyone would then
offload. Compute is not the only reason to centralize: pushing $15$\,GB of
weights to every robot for every version is an operational cost a task-agnostic
$0.14$\,GB world model does not incur; frontier weights are often not
distributed at all; centralized execution is what lets fleet experience be
aggregated; and one server amortized over $N$
robots is the deployment in which uplink is genuinely contended. We did not
evaluate the multi-robot case and consider it the most interesting missing
experiment.

\paragraph{What the head is, and is not.} $\Phi$ estimates a quantity defined
by a specific policy, so it inherits that policy's blind spots: if the policy
is indifferent to something it should attend to, \sys will happily not transmit
it. The head is also trained, which sits uneasily beside our claim that the
world model's quality does not matter --- it does not matter for
\emph{Figure~\ref{fig:m2}}, which is measurement, but it does matter for the
system. A head miscalibrated for a new embodiment degrades toward the
world-model-latent trigger, which \S\ref{sec:eval-ablation} measures at
$-2.5$\pp rather than at collapse, so the failure is graceful; we would still
want an online confidence check before deploying on hardware we have not
calibrated.

\paragraph{What we did not evaluate.} A compromised client can suppress
transmission by reporting low $\hat{\delta}_t$; server-side plausibility checks
against the heartbeat's own divergence estimate would catch it and are
unimplemented. Every closed-loop result is a $7$-DoF arm.
And our physical experiment shapes uplink capacity rather than replaying a live
cellular link, so it validates the response to scarcity and not the interaction
with real radio dynamics.

%%%%%%%%%%%%%%%%%%%%%%%%%%%%%%%%%%%%%%%%%%%%%%%%%%%%%%%%%%%%%%%%%%%%%%%%%%%%%%%
\section{Conclusion}
\label{sec:conclusion}

Predictive transmission has been asking the same question for twenty years: how
wrong is the receiver's reconstruction? That question is right when the receiver
renders what it receives, and wrong when the receiver is a policy that acts on
what it receives and thereby chooses what it receives next. We measured the
difference --- fidelity ranks decision impact at $\rho_s = 0.31$ and
misclassifies $44\%$ of frames in both directions --- and built the transmitter
that asks the other question: a shared action-conditioned world model, a
distilled robot-side estimate of how much a frame would move the policy's
action, a bandwidth-adaptive threshold, and the resynchronization protocol that
closed-loop prediction needs and event-triggered sampling never had. At the same
world model and the same task success, that change of criterion alone is worth
$59\%$ of the uplink, $16\times$ against the deployed status quo and $75\times$
composed with spatial compression. The general claim is shorter: when the
consumer of a stream is a decision-maker, the transmitter should be measuring
decisions, not pixels --- and it can, cheaply, without ever running the
decision-maker.

%%%%%%%%%%%%%%%%%%%%%%%%%%%%%%%%%%%%%%%%%%%%%%%%%%%%%%%%%%%%%%%%%%%%%%%%%%%%%%%
%% References do not count against the 12-page limit.
\bibliographystyle{ACM-Reference-Format}
\bibliography{references}

%%%%%%%%%%%%%%%%%%%%%%%%%%%%%%%%%%%%%%%%%%%%%%%%%%%%%%%%%%%%%%%%%%%%%%%%%%%%%%%
%% Appendices go AFTER the bibliography and do not count against the 12 pages.
% \appendix
% \section{Additional Results}

\end{document}
